t%%%%%%%%%%%%%%%%%%%%%%%%%%%%%%%%%%%%%%%%%%%%%%%%%%%%%%%%%%%%
% 14.13 PROBABILITY DISTRIBUTIONS
% The Composition of Advanced Mathematics
%%%%%%%%%%%%%%%%%%%%%%%%%%%%%%%%%%%%%%%%%%%%%%%%%%%%%%%%%%%%

\section{Probability Distributions}
\label{sec:probability-distributions}

\prerequisite{
Elementary probability, random variables, expectation,
variance, combinatorics, integrals, exponential functions,
Gamma functions, and basic statistical notation.
}

\learningobjectives{
By the end of this reference section, the reader should be able to:
\begin{itemize}
    \item distinguish discrete and continuous random variables;
    \item recognize probability mass, density, and distribution functions;
    \item identify supports and parameter restrictions;
    \item compute expectation and variance from standard distributions;
    \item recognize Bernoulli, binomial, geometric, negative-binomial,
          hypergeometric, and Poisson distributions;
    \item recognize discrete-uniform distributions;
    \item recognize continuous-uniform, exponential, Gamma, Beta,
          normal, lognormal, Weibull, Cauchy, chi-square,
          Student's \(t\), and \(F\) distributions;
    \item understand common parameterization differences;
    \item use standardization for normal variables;
    \item recognize relationships among distributions;
    \item recognize memorylessness;
    \item identify moment-generating and characteristic functions;
    \item distinguish population parameters from realized observations;
    \item select a plausible standard model from a problem description;
    \item avoid common parameter and support errors.
\end{itemize}
}

%%%%%%%%%%%%%%%%%%%%%%%%%%%%%%%%%%%%%%%%%%%%%%%%%%%%%%%%%%%%
\subsection{Random Variables}
\label{subsec:random-variable-reference}
%%%%%%%%%%%%%%%%%%%%%%%%%%%%%%%%%%%%%%%%%%%%%%%%%%%%%%%%%%%%

A random variable assigns a numerical value to each outcome of a random
experiment.

A random variable is commonly classified as:

\[
\boxed{
\text{discrete}
}
\]

or

\[
\boxed{
\text{continuous}.
}
\]

%%%%%%%%%%%%%%%%%%%%%%%%%%%%%%%%%%%%%%%%%%%%%%%%%%%%%%%%%%%%
\subsection{Discrete Probability Mass Function}
\label{subsec:pmf-reference}
%%%%%%%%%%%%%%%%%%%%%%%%%%%%%%%%%%%%%%%%%%%%%%%%%%%%%%%%%%%%

For a discrete random variable \(X\), the probability mass function is

\[
\boxed{
p_X(x)
=
\Pr(X=x).
}
\]

It satisfies

\[
\boxed{
p_X(x)\geq0
}
\]

and

\[
\boxed{
\sum_x p_X(x)=1.
}
\]

%%%%%%%%%%%%%%%%%%%%%%%%%%%%%%%%%%%%%%%%%%%%%%%%%%%%%%%%%%%%
\subsection{Continuous Probability Density Function}
\label{subsec:pdf-reference}
%%%%%%%%%%%%%%%%%%%%%%%%%%%%%%%%%%%%%%%%%%%%%%%%%%%%%%%%%%%%

For a continuous random variable \(X\), a probability density function
\(f_X(x)\) satisfies

\[
\boxed{
f_X(x)\geq0
}
\]

and

\[
\boxed{
\int_{-\infty}^{\infty}
f_X(x)\,dx
=
1.
}
\]

Probabilities are computed by integration:

\[
\boxed{
\Pr(a\leq X\leq b)
=
\int_a^b
f_X(x)\,dx.
}
\]

For a continuous random variable,

\[
\boxed{
\Pr(X=x)=0.
}
\]

%%%%%%%%%%%%%%%%%%%%%%%%%%%%%%%%%%%%%%%%%%%%%%%%%%%%%%%%%%%%
\subsection{Cumulative Distribution Function}
\label{subsec:cdf-reference}
%%%%%%%%%%%%%%%%%%%%%%%%%%%%%%%%%%%%%%%%%%%%%%%%%%%%%%%%%%%%

For any random variable,

\[
\boxed{
F_X(x)
=
\Pr(X\leq x).
}
\]

The cumulative distribution function, or CDF, satisfies

\[
\boxed{
0\leq F_X(x)\leq1,
}
\]

is nondecreasing, and obeys

\[
\boxed{
\lim_{x\to-\infty}F_X(x)=0,
}
\]

\[
\boxed{
\lim_{x\to\infty}F_X(x)=1.
}
\]

%%%%%%%%%%%%%%%%%%%%%%%%%%%%%%%%%%%%%%%%%%%%%%%%%%%%%%%%%%%%
\subsection{Expectation}
\label{subsec:expectation-distribution-reference}
%%%%%%%%%%%%%%%%%%%%%%%%%%%%%%%%%%%%%%%%%%%%%%%%%%%%%%%%%%%%

For a discrete random variable,

\[
\boxed{
\mathbb{E}[X]
=
\sum_x
x\,p_X(x).
}
\]

For a continuous random variable,

\[
\boxed{
\mathbb{E}[X]
=
\int_{-\infty}^{\infty}
x f_X(x)\,dx.
}
\]

%%%%%%%%%%%%%%%%%%%%%%%%%%%%%%%%%%%%%%%%%%%%%%%%%%%%%%%%%%%%
\subsection{Variance}
\label{subsec:variance-distribution-reference}
%%%%%%%%%%%%%%%%%%%%%%%%%%%%%%%%%%%%%%%%%%%%%%%%%%%%%%%%%%%%

The variance is

\[
\boxed{
\operatorname{Var}(X)
=
\mathbb{E}
\left[
(X-\mu)^2
\right],
}
\]

where

\[
\mu=\mathbb{E}[X].
\]

Equivalent form:

\[
\boxed{
\operatorname{Var}(X)
=
\mathbb{E}[X^2]
-
(\mathbb{E}[X])^2.
}
\]

Standard deviation is

\[
\boxed{
\sigma
=
\sqrt{
\operatorname{Var}(X)
}.
}
\]

%%%%%%%%%%%%%%%%%%%%%%%%%%%%%%%%%%%%%%%%%%%%%%%%%%%%%%%%%%%%
\subsection{Bernoulli Distribution}
\label{subsec:bernoulli-distribution-reference}
%%%%%%%%%%%%%%%%%%%%%%%%%%%%%%%%%%%%%%%%%%%%%%%%%%%%%%%%%%%%

Notation:

\[
\boxed{
X\sim\operatorname{Bernoulli}(p)
}
\]

with

\[
0\leq p\leq1.
\]

Support:

\[
\boxed{
x\in\{0,1\}.
}
\]

PMF:

\[
\boxed{
\Pr(X=x)
=
p^x(1-p)^{1-x}.
}
\]

Mean:

\[
\boxed{
\mathbb{E}[X]=p.
}
\]

Variance:

\[
\boxed{
\operatorname{Var}(X)
=
p(1-p).
}
\]

%%%%%%%%%%%%%%%%%%%%%%%%%%%%%%%%%%%%%%%%%%%%%%%%%%%%%%%%%%%%
\subsection{Bernoulli Interpretation}
\label{subsec:bernoulli-interpretation}
%%%%%%%%%%%%%%%%%%%%%%%%%%%%%%%%%%%%%%%%%%%%%%%%%%%%%%%%%%%%

A Bernoulli random variable models one trial with two outcomes.

Common coding is

\[
\boxed{
X=
\begin{cases}
1, & \text{success},\\
0, & \text{failure}.
\end{cases}
}
\]

%%%%%%%%%%%%%%%%%%%%%%%%%%%%%%%%%%%%%%%%%%%%%%%%%%%%%%%%%%%%
\subsection{Binomial Distribution}
\label{subsec:binomial-distribution-reference}
%%%%%%%%%%%%%%%%%%%%%%%%%%%%%%%%%%%%%%%%%%%%%%%%%%%%%%%%%%%%

Notation:

\[
\boxed{
X\sim\operatorname{Binomial}(n,p).
}
\]

The distribution models the number of successes in \(n\) independent
Bernoulli trials with constant success probability \(p\).

Support:

\[
\boxed{
x=0,1,\ldots,n.
}
\]

PMF:

\[
\boxed{
\Pr(X=x)
=
\binom{n}{x}
p^x(1-p)^{n-x}.
}
\]

Mean:

\[
\boxed{
\mathbb{E}[X]
=
np.
}
\]

Variance:

\[
\boxed{
\operatorname{Var}(X)
=
np(1-p).
}
\]

%%%%%%%%%%%%%%%%%%%%%%%%%%%%%%%%%%%%%%%%%%%%%%%%%%%%%%%%%%%%
\subsection{Worked Example: Binomial Probability}
\label{subsec:worked-binomial-probability}
%%%%%%%%%%%%%%%%%%%%%%%%%%%%%%%%%%%%%%%%%%%%%%%%%%%%%%%%%%%%

\begin{workedexamplebox}{Exactly Three Successes}

Suppose

\[
X\sim\operatorname{Binomial}(5,0.4).
\]

Find

\[
\Pr(X=3).
\]

Use

\[
\Pr(X=3)
=
\binom53
(0.4)^3
(0.6)^2.
\]

Thus

\begin{answerbox}
\[
\boxed{
\Pr(X=3)
=
10(0.4)^3(0.6)^2
=
0.2304.
}
\]
\end{answerbox}

\end{workedexamplebox}

%%%%%%%%%%%%%%%%%%%%%%%%%%%%%%%%%%%%%%%%%%%%%%%%%%%%%%%%%%%%
\subsection{Geometric Distribution}
\label{subsec:geometric-distribution-reference}
%%%%%%%%%%%%%%%%%%%%%%%%%%%%%%%%%%%%%%%%%%%%%%%%%%%%%%%%%%%%

Using the convention that \(X\) counts the trial on which the first success
occurs,

\[
\boxed{
X\sim\operatorname{Geometric}(p).
}
\]

Support:

\[
\boxed{
x=1,2,3,\ldots.
}
\]

PMF:

\[
\boxed{
\Pr(X=x)
=
(1-p)^{x-1}p.
}
\]

Mean:

\[
\boxed{
\mathbb{E}[X]
=
\frac1p.
}
\]

Variance:

\[
\boxed{
\operatorname{Var}(X)
=
\frac{1-p}{p^2}.
}
\]

%%%%%%%%%%%%%%%%%%%%%%%%%%%%%%%%%%%%%%%%%%%%%%%%%%%%%%%%%%%%
\subsection{Geometric Parameterization Warning}
\label{subsec:geometric-parameterization-warning}
%%%%%%%%%%%%%%%%%%%%%%%%%%%%%%%%%%%%%%%%%%%%%%%%%%%%%%%%%%%%

Some texts define the geometric random variable as the number of failures
before the first success.

Under that convention, support begins at \(0\).

Always check the definition before using formulas.

%%%%%%%%%%%%%%%%%%%%%%%%%%%%%%%%%%%%%%%%%%%%%%%%%%%%%%%%%%%%
\subsection{Memoryless Property of the Geometric Distribution}
\label{subsec:geometric-memoryless}
%%%%%%%%%%%%%%%%%%%%%%%%%%%%%%%%%%%%%%%%%%%%%%%%%%%%%%%%%%%%

For nonnegative integers \(m,n\),

\[
\boxed{
\Pr(X>m+n\mid X>m)
=
\Pr(X>n).
}
\]

The geometric distribution is the standard discrete memoryless distribution.

%%%%%%%%%%%%%%%%%%%%%%%%%%%%%%%%%%%%%%%%%%%%%%%%%%%%%%%%%%%%
\subsection{Negative Binomial Distribution}
\label{subsec:negative-binomial-reference}
%%%%%%%%%%%%%%%%%%%%%%%%%%%%%%%%%%%%%%%%%%%%%%%%%%%%%%%%%%%%

One common convention lets \(X\) count the trial on which the \(r\)-th
success occurs.

Then

\[
\boxed{
X\sim\operatorname{NegBin}(r,p),
}
\]

with support

\[
\boxed{
x=r,r+1,\ldots.
}
\]

PMF:

\[
\boxed{
\Pr(X=x)
=
\binom{x-1}{r-1}
p^r(1-p)^{x-r}.
}
\]

Mean:

\[
\boxed{
\mathbb{E}[X]
=
\frac{r}{p}.
}
\]

Variance:

\[
\boxed{
\operatorname{Var}(X)
=
\frac{
r(1-p)
}{
p^2
}.
}
\]

%%%%%%%%%%%%%%%%%%%%%%%%%%%%%%%%%%%%%%%%%%%%%%%%%%%%%%%%%%%%
\subsection{Hypergeometric Distribution}
\label{subsec:hypergeometric-distribution-reference}
%%%%%%%%%%%%%%%%%%%%%%%%%%%%%%%%%%%%%%%%%%%%%%%%%%%%%%%%%%%%

Suppose a population contains:

\[
N
=
\text{total objects},
\]

\[
K
=
\text{success-type objects},
\]

and a sample of size \(n\) is drawn without replacement.

Then

\[
\boxed{
X\sim\operatorname{Hypergeometric}(N,K,n).
}
\]

PMF:

\[
\boxed{
\Pr(X=x)
=
\frac{
\binom{K}{x}
\binom{N-K}{n-x}
}{
\binom{N}{n}
}.
}
\]

Mean:

\[
\boxed{
\mathbb{E}[X]
=
n\frac{K}{N}.
}
\]

Variance:

\[
\boxed{
\operatorname{Var}(X)
=
n
\frac{K}{N}
\left(
1-\frac{K}{N}
\right)
\frac{
N-n
}{
N-1
}.
}
\]

%%%%%%%%%%%%%%%%%%%%%%%%%%%%%%%%%%%%%%%%%%%%%%%%%%%%%%%%%%%%
\subsection{Binomial Versus Hypergeometric}
\label{subsec:binomial-vs-hypergeometric}
%%%%%%%%%%%%%%%%%%%%%%%%%%%%%%%%%%%%%%%%%%%%%%%%%%%%%%%%%%%%

Use the binomial distribution for independent trials with constant success
probability.

Use the hypergeometric distribution for sampling without replacement from a
finite population.

%%%%%%%%%%%%%%%%%%%%%%%%%%%%%%%%%%%%%%%%%%%%%%%%%%%%%%%%%%%%
\subsection{Poisson Distribution}
\label{subsec:poisson-distribution-reference}
%%%%%%%%%%%%%%%%%%%%%%%%%%%%%%%%%%%%%%%%%%%%%%%%%%%%%%%%%%%%

Notation:

\[
\boxed{
X\sim\operatorname{Poisson}(\lambda),
\qquad
\lambda>0.
}
\]

Support:

\[
\boxed{
x=0,1,2,\ldots.
}
\]

PMF:

\[
\boxed{
\Pr(X=x)
=
e^{-\lambda}
\frac{
\lambda^x
}{
x!
}.
}
\]

Mean:

\[
\boxed{
\mathbb{E}[X]
=
\lambda.
}
\]

Variance:

\[
\boxed{
\operatorname{Var}(X)
=
\lambda.
}
\]

%%%%%%%%%%%%%%%%%%%%%%%%%%%%%%%%%%%%%%%%%%%%%%%%%%%%%%%%%%%%
\subsection{Poisson Interpretation}
\label{subsec:poisson-interpretation}
%%%%%%%%%%%%%%%%%%%%%%%%%%%%%%%%%%%%%%%%%%%%%%%%%%%%%%%%%%%%

The Poisson distribution commonly models event counts in a fixed interval
when events occur according to a homogeneous Poisson-process model.

Examples include counts over:

\[
\boxed{
\text{time},
\quad
\text{area},
\quad
\text{length},
\quad
\text{volume}.
}
\]

%%%%%%%%%%%%%%%%%%%%%%%%%%%%%%%%%%%%%%%%%%%%%%%%%%%%%%%%%%%%
\subsection{Poisson Approximation to Binomial}
\label{subsec:poisson-binomial-approximation}
%%%%%%%%%%%%%%%%%%%%%%%%%%%%%%%%%%%%%%%%%%%%%%%%%%%%%%%%%%%%

For large \(n\), small \(p\), and moderate

\[
\lambda=np,
\]

a binomial variable can often be approximated by

\[
\boxed{
X
\approx
\operatorname{Poisson}(np).
}
\]

The quality depends on the specific parameter regime.

%%%%%%%%%%%%%%%%%%%%%%%%%%%%%%%%%%%%%%%%%%%%%%%%%%%%%%%%%%%%
\subsection{Discrete Uniform Distribution}
\label{subsec:discrete-uniform-reference}
%%%%%%%%%%%%%%%%%%%%%%%%%%%%%%%%%%%%%%%%%%%%%%%%%%%%%%%%%%%%

If \(X\) is equally likely to take values

\[
1,2,\ldots,n,
\]

then

\[
\boxed{
\Pr(X=x)
=
\frac1n.
}
\]

Mean:

\[
\boxed{
\mathbb{E}[X]
=
\frac{n+1}{2}.
}
\]

Variance:

\[
\boxed{
\operatorname{Var}(X)
=
\frac{
n^2-1
}{
12
}.
}
\]

%%%%%%%%%%%%%%%%%%%%%%%%%%%%%%%%%%%%%%%%%%%%%%%%%%%%%%%%%%%%
\subsection{Continuous Uniform Distribution}
\label{subsec:continuous-uniform-reference}
%%%%%%%%%%%%%%%%%%%%%%%%%%%%%%%%%%%%%%%%%%%%%%%%%%%%%%%%%%%%

Notation:

\[
\boxed{
X\sim\operatorname{Uniform}(a,b),
\qquad
a<b.
}
\]

Density:

\[
\boxed{
f(x)
=
\frac{
1
}{
b-a
},
\qquad
a\leq x\leq b.
}
\]

Mean:

\[
\boxed{
\mathbb{E}[X]
=
\frac{a+b}{2}.
}
\]

Variance:

\[
\boxed{
\operatorname{Var}(X)
=
\frac{
(b-a)^2
}{
12
}.
}
\]

%%%%%%%%%%%%%%%%%%%%%%%%%%%%%%%%%%%%%%%%%%%%%%%%%%%%%%%%%%%%
\subsection{Exponential Distribution}
\label{subsec:exponential-distribution-reference}
%%%%%%%%%%%%%%%%%%%%%%%%%%%%%%%%%%%%%%%%%%%%%%%%%%%%%%%%%%%%

Using rate parameter \(\lambda>0\),

\[
\boxed{
X\sim\operatorname{Exponential}(\lambda).
}
\]

Support:

\[
\boxed{
x\geq0.
}
\]

Density:

\[
\boxed{
f(x)
=
\lambda e^{-\lambda x}.
}
\]

CDF:

\[
\boxed{
F(x)
=
1-e^{-\lambda x},
\qquad
x\geq0.
}
\]

Mean:

\[
\boxed{
\mathbb{E}[X]
=
\frac1\lambda.
}
\]

Variance:

\[
\boxed{
\operatorname{Var}(X)
=
\frac1{\lambda^2}.
}
\]

%%%%%%%%%%%%%%%%%%%%%%%%%%%%%%%%%%%%%%%%%%%%%%%%%%%%%%%%%%%%
\subsection{Exponential Memoryless Property}
\label{subsec:exponential-memoryless}
%%%%%%%%%%%%%%%%%%%%%%%%%%%%%%%%%%%%%%%%%%%%%%%%%%%%%%%%%%%%

For \(s,t\geq0\),

\[
\boxed{
\Pr(X>s+t\mid X>s)
=
\Pr(X>t).
}
\]

The exponential distribution is the standard continuous memoryless
distribution.

%%%%%%%%%%%%%%%%%%%%%%%%%%%%%%%%%%%%%%%%%%%%%%%%%%%%%%%%%%%%
\subsection{Gamma Distribution}
\label{subsec:gamma-distribution-reference}
%%%%%%%%%%%%%%%%%%%%%%%%%%%%%%%%%%%%%%%%%%%%%%%%%%%%%%%%%%%%

Using shape \(\alpha>0\) and rate \(\beta>0\),

\[
\boxed{
X\sim\operatorname{Gamma}(\alpha,\beta).
}
\]

Density:

\[
\boxed{
f(x)
=
\frac{
\beta^\alpha
}{
\Gamma(\alpha)
}
x^{\alpha-1}
e^{-\beta x},
\qquad
x>0.
}
\]

Mean:

\[
\boxed{
\mathbb{E}[X]
=
\frac{\alpha}{\beta}.
}
\]

Variance:

\[
\boxed{
\operatorname{Var}(X)
=
\frac{\alpha}{\beta^2}.
}
\]

%%%%%%%%%%%%%%%%%%%%%%%%%%%%%%%%%%%%%%%%%%%%%%%%%%%%%%%%%%%%
\subsection{Gamma Scale Parameterization}
\label{subsec:gamma-scale-parameterization}
%%%%%%%%%%%%%%%%%%%%%%%%%%%%%%%%%%%%%%%%%%%%%%%%%%%%%%%%%%%%

Some texts use scale

\[
\theta
=
\frac1\beta.
\]

Then

\[
\boxed{
\mathbb{E}[X]
=
\alpha\theta,
}
\]

and

\[
\boxed{
\operatorname{Var}(X)
=
\alpha\theta^2.
}
\]

Always determine whether the second parameter is a rate or scale.

%%%%%%%%%%%%%%%%%%%%%%%%%%%%%%%%%%%%%%%%%%%%%%%%%%%%%%%%%%%%
\subsection{Erlang Distribution}
\label{subsec:erlang-distribution-reference}
%%%%%%%%%%%%%%%%%%%%%%%%%%%%%%%%%%%%%%%%%%%%%%%%%%%%%%%%%%%%

The Erlang distribution is a Gamma distribution with integer shape parameter.

If

\[
\alpha=k\in\mathbb{N},
\]

then

\[
\boxed{
X
\sim
\operatorname{Erlang}(k,\lambda).
}
\]

It commonly models the waiting time until the \(k\)-th event in a Poisson
process.

%%%%%%%%%%%%%%%%%%%%%%%%%%%%%%%%%%%%%%%%%%%%%%%%%%%%%%%%%%%%
\subsection{Beta Distribution}
\label{subsec:beta-distribution-reference}
%%%%%%%%%%%%%%%%%%%%%%%%%%%%%%%%%%%%%%%%%%%%%%%%%%%%%%%%%%%%

For

\[
\alpha>0,
\qquad
\beta>0,
\]

\[
\boxed{
X\sim\operatorname{Beta}(\alpha,\beta).
}
\]

Support:

\[
\boxed{
0<x<1.
}
\]

Density:

\[
\boxed{
f(x)
=
\frac{
x^{\alpha-1}
(1-x)^{\beta-1}
}{
B(\alpha,\beta)
}.
}
\]

Mean:

\[
\boxed{
\mathbb{E}[X]
=
\frac{
\alpha
}{
\alpha+\beta
}.
}
\]

Variance:

\[
\boxed{
\operatorname{Var}(X)
=
\frac{
\alpha\beta
}{
(\alpha+\beta)^2
(\alpha+\beta+1)
}.
}
\]

%%%%%%%%%%%%%%%%%%%%%%%%%%%%%%%%%%%%%%%%%%%%%%%%%%%%%%%%%%%%
\subsection{Normal Distribution}
\label{subsec:normal-distribution-reference}
%%%%%%%%%%%%%%%%%%%%%%%%%%%%%%%%%%%%%%%%%%%%%%%%%%%%%%%%%%%%

Notation:

\[
\boxed{
X\sim N(\mu,\sigma^2),
\qquad
\sigma>0.
}
\]

Density:

\[
\boxed{
f(x)
=
\frac{
1
}{
\sigma\sqrt{2\pi}
}
\exp
\left[
-\frac{
(x-\mu)^2
}{
2\sigma^2
}
\right].
}
\]

Mean:

\[
\boxed{
\mathbb{E}[X]=\mu.
}
\]

Variance:

\[
\boxed{
\operatorname{Var}(X)=\sigma^2.
}
\]

%%%%%%%%%%%%%%%%%%%%%%%%%%%%%%%%%%%%%%%%%%%%%%%%%%%%%%%%%%%%
\subsection{Standard Normal Distribution}
\label{subsec:standard-normal-reference}
%%%%%%%%%%%%%%%%%%%%%%%%%%%%%%%%%%%%%%%%%%%%%%%%%%%%%%%%%%%%

The standard normal variable is

\[
\boxed{
Z\sim N(0,1).
}
\]

Density:

\[
\boxed{
\phi(z)
=
\frac{
1
}{
\sqrt{2\pi}
}
e^{-z^2/2}.
}
\]

CDF:

\[
\boxed{
\Phi(z)
=
\Pr(Z\leq z).
}
\]

%%%%%%%%%%%%%%%%%%%%%%%%%%%%%%%%%%%%%%%%%%%%%%%%%%%%%%%%%%%%
\subsection{Normal Standardization}
\label{subsec:normal-standardization-reference}
%%%%%%%%%%%%%%%%%%%%%%%%%%%%%%%%%%%%%%%%%%%%%%%%%%%%%%%%%%%%

If

\[
X\sim N(\mu,\sigma^2),
\]

then

\[
\boxed{
Z
=
\frac{
X-\mu
}{
\sigma
}
\sim N(0,1).
}
\]

Thus

\[
\boxed{
\Pr(X\leq x)
=
\Phi
\left(
\frac{
x-\mu
}{
\sigma
}
\right).
}
\]

%%%%%%%%%%%%%%%%%%%%%%%%%%%%%%%%%%%%%%%%%%%%%%%%%%%%%%%%%%%%
\subsection{Worked Example: Normal Standardization}
\label{subsec:worked-normal-standardization}
%%%%%%%%%%%%%%%%%%%%%%%%%%%%%%%%%%%%%%%%%%%%%%%%%%%%%%%%%%%%

\begin{workedexamplebox}{Convert to a \(z\)-Score}

Suppose

\[
X\sim N(100,15^2).
\]

For

\[
x=130,
\]

the standardized value is

\[
z
=
\frac{
130-100
}{
15
}
=
2.
\]

Therefore,

\begin{answerbox}
\[
\boxed{
\Pr(X\leq130)
=
\Phi(2).
}
\]
\end{answerbox}

\end{workedexamplebox}

%%%%%%%%%%%%%%%%%%%%%%%%%%%%%%%%%%%%%%%%%%%%%%%%%%%%%%%%%%%%
\subsection{Normal Empirical Rule}
\label{subsec:normal-empirical-rule}
%%%%%%%%%%%%%%%%%%%%%%%%%%%%%%%%%%%%%%%%%%%%%%%%%%%%%%%%%%%%

For a normal distribution, approximately:

\[
\boxed{
68\%
}
\]

of observations lie within one standard deviation of the mean,

\[
\boxed{
95\%
}
\]

within two standard deviations, and

\[
\boxed{
99.7\%
}
\]

within three standard deviations.

This is an approximation, not an exact identity.

%%%%%%%%%%%%%%%%%%%%%%%%%%%%%%%%%%%%%%%%%%%%%%%%%%%%%%%%%%%%
\subsection{Lognormal Distribution}
\label{subsec:lognormal-distribution-reference}
%%%%%%%%%%%%%%%%%%%%%%%%%%%%%%%%%%%%%%%%%%%%%%%%%%%%%%%%%%%%

A positive random variable \(X\) is lognormal if

\[
\boxed{
\ln X
\sim
N(\mu,\sigma^2).
}
\]

Notation:

\[
\boxed{
X
\sim
\operatorname{Lognormal}(\mu,\sigma^2).
}
\]

Mean:

\[
\boxed{
\mathbb{E}[X]
=
e^{\mu+\sigma^2/2}.
}
\]

Variance:

\[
\boxed{
\operatorname{Var}(X)
=
\left(
e^{\sigma^2}-1
\right)
e^{2\mu+\sigma^2}.
}
\]

%%%%%%%%%%%%%%%%%%%%%%%%%%%%%%%%%%%%%%%%%%%%%%%%%%%%%%%%%%%%
\subsection{Weibull Distribution}
\label{subsec:weibull-distribution-reference}
%%%%%%%%%%%%%%%%%%%%%%%%%%%%%%%%%%%%%%%%%%%%%%%%%%%%%%%%%%%%

Using shape \(k>0\) and scale \(\lambda>0\),

\[
\boxed{
X\sim\operatorname{Weibull}(k,\lambda).
}
\]

Density:

\[
\boxed{
f(x)
=
\frac{k}{\lambda}
\left(
\frac{x}{\lambda}
\right)^{k-1}
\exp
\left[
-
\left(
\frac{x}{\lambda}
\right)^k
\right],
\qquad
x\geq0.
}
\]

CDF:

\[
\boxed{
F(x)
=
1-
\exp
\left[
-
\left(
\frac{x}{\lambda}
\right)^k
\right].
}
\]

Mean:

\[
\boxed{
\mathbb{E}[X]
=
\lambda
\Gamma
\left(
1+\frac1k
\right).
}
\]

%%%%%%%%%%%%%%%%%%%%%%%%%%%%%%%%%%%%%%%%%%%%%%%%%%%%%%%%%%%%
\subsection{Cauchy Distribution}
\label{subsec:cauchy-distribution-reference}
%%%%%%%%%%%%%%%%%%%%%%%%%%%%%%%%%%%%%%%%%%%%%%%%%%%%%%%%%%%%

Using location \(x_0\) and scale \(\gamma>0\),

\[
\boxed{
X\sim\operatorname{Cauchy}(x_0,\gamma).
}
\]

Density:

\[
\boxed{
f(x)
=
\frac{
1
}{
\pi\gamma
\left[
1+
\left(
\frac{x-x_0}{\gamma}
\right)^2
\right]
}.
}
\]

The Cauchy distribution does not have a finite ordinary mean or variance.

%%%%%%%%%%%%%%%%%%%%%%%%%%%%%%%%%%%%%%%%%%%%%%%%%%%%%%%%%%%%
\subsection{Chi-Square Distribution}
\label{subsec:chi-square-distribution-reference}
%%%%%%%%%%%%%%%%%%%%%%%%%%%%%%%%%%%%%%%%%%%%%%%%%%%%%%%%%%%%

Notation:

\[
\boxed{
X\sim\chi_\nu^2,
}
\]

where \(\nu>0\) is the degrees of freedom.

A chi-square distribution is a Gamma distribution with

\[
\boxed{
\alpha=\frac{\nu}{2},
\qquad
\beta=\frac12
}
\]

using the rate parameterization.

Mean:

\[
\boxed{
\mathbb{E}[X]
=
\nu.
}
\]

Variance:

\[
\boxed{
\operatorname{Var}(X)
=
2\nu.
}
\]

%%%%%%%%%%%%%%%%%%%%%%%%%%%%%%%%%%%%%%%%%%%%%%%%%%%%%%%%%%%%
\subsection{Chi-Square Construction}
\label{subsec:chi-square-construction}
%%%%%%%%%%%%%%%%%%%%%%%%%%%%%%%%%%%%%%%%%%%%%%%%%%%%%%%%%%%%

If

\[
Z_1,\ldots,Z_\nu
\]

are independent standard normal variables, then

\[
\boxed{
Z_1^2+\cdots+Z_\nu^2
\sim
\chi_\nu^2.
}
\]

%%%%%%%%%%%%%%%%%%%%%%%%%%%%%%%%%%%%%%%%%%%%%%%%%%%%%%%%%%%%
\subsection{Student's \(t\) Distribution}
\label{subsec:student-t-distribution-reference}
%%%%%%%%%%%%%%%%%%%%%%%%%%%%%%%%%%%%%%%%%%%%%%%%%%%%%%%%%%%%

Notation:

\[
\boxed{
T\sim t_\nu.
}
\]

If

\[
Z\sim N(0,1),
\]

\[
V\sim\chi_\nu^2,
\]

and \(Z\) and \(V\) are independent, then

\[
\boxed{
T
=
\frac{
Z
}{
\sqrt{V/\nu}
}
\sim
t_\nu.
}
\]

The distribution is symmetric about \(0\).

For \(\nu>1\),

\[
\boxed{
\mathbb{E}[T]=0.
}
\]

For \(\nu>2\),

\[
\boxed{
\operatorname{Var}(T)
=
\frac{
\nu
}{
\nu-2
}.
}
\]

%%%%%%%%%%%%%%%%%%%%%%%%%%%%%%%%%%%%%%%%%%%%%%%%%%%%%%%%%%%%
\subsection{Student's \(t\) and Normal Limit}
\label{subsec:t-normal-limit-reference}
%%%%%%%%%%%%%%%%%%%%%%%%%%%%%%%%%%%%%%%%%%%%%%%%%%%%%%%%%%%%

As

\[
\nu\to\infty,
\]

the Student's \(t\) distribution approaches the standard normal:

\[
\boxed{
t_\nu
\Longrightarrow
N(0,1).
}
\]

%%%%%%%%%%%%%%%%%%%%%%%%%%%%%%%%%%%%%%%%%%%%%%%%%%%%%%%%%%%%
\subsection{\(F\) Distribution}
\label{subsec:f-distribution-reference}
%%%%%%%%%%%%%%%%%%%%%%%%%%%%%%%%%%%%%%%%%%%%%%%%%%%%%%%%%%%%

If

\[
U_1\sim\chi_{\nu_1}^2,
\qquad
U_2\sim\chi_{\nu_2}^2
\]

are independent, then

\[
\boxed{
F
=
\frac{
U_1/\nu_1
}{
U_2/\nu_2
}
\sim
F_{\nu_1,\nu_2}.
}
\]

The \(F\) distribution is supported on

\[
\boxed{
x>0.
}
\]

It is central in variance comparisons and analysis of variance.

%%%%%%%%%%%%%%%%%%%%%%%%%%%%%%%%%%%%%%%%%%%%%%%%%%%%%%%%%%%%
\subsection{Relationships Among Common Distributions}
\label{subsec:distribution-relationships-reference}
%%%%%%%%%%%%%%%%%%%%%%%%%%%%%%%%%%%%%%%%%%%%%%%%%%%%%%%%%%%%

Important relationships include:

\[
\boxed{
\operatorname{Bernoulli}(p)
=
\operatorname{Binomial}(1,p).
}
\]

The sum of independent Bernoulli variables with common \(p\) is binomial.

The sum of independent Poisson variables is Poisson:

\[
\boxed{
X_i\sim\operatorname{Poisson}(\lambda_i)
\Longrightarrow
\sum_iX_i
\sim
\operatorname{Poisson}
\left(
\sum_i\lambda_i
\right).
}
\]

The exponential distribution is Gamma with shape \(1\).

The chi-square distribution is a special Gamma distribution.

%%%%%%%%%%%%%%%%%%%%%%%%%%%%%%%%%%%%%%%%%%%%%%%%%%%%%%%%%%%%
\subsection{Normal Sums}
\label{subsec:normal-sums-reference}
%%%%%%%%%%%%%%%%%%%%%%%%%%%%%%%%%%%%%%%%%%%%%%%%%%%%%%%%%%%%

If independent

\[
X_i
\sim
N(\mu_i,\sigma_i^2),
\]

then

\[
\boxed{
\sum_i X_i
\sim
N
\left(
\sum_i\mu_i,
\sum_i\sigma_i^2
\right).
}
\]

More generally, linear combinations of jointly normal variables are normal.

%%%%%%%%%%%%%%%%%%%%%%%%%%%%%%%%%%%%%%%%%%%%%%%%%%%%%%%%%%%%
\subsection{Moment-Generating Function}
\label{subsec:mgf-distribution-reference}
%%%%%%%%%%%%%%%%%%%%%%%%%%%%%%%%%%%%%%%%%%%%%%%%%%%%%%%%%%%%

The moment-generating function is

\[
\boxed{
M_X(t)
=
\mathbb{E}[e^{tX}]
}
\]

when it exists near \(t=0\).

Moments can be obtained through

\[
\boxed{
M_X^{(n)}(0)
=
\mathbb{E}[X^n].
}
\]

%%%%%%%%%%%%%%%%%%%%%%%%%%%%%%%%%%%%%%%%%%%%%%%%%%%%%%%%%%%%
\subsection{Characteristic Function}
\label{subsec:characteristic-function-distribution-reference}
%%%%%%%%%%%%%%%%%%%%%%%%%%%%%%%%%%%%%%%%%%%%%%%%%%%%%%%%%%%%

The characteristic function is

\[
\boxed{
\varphi_X(t)
=
\mathbb{E}[e^{itX}].
}
\]

Unlike the moment-generating function, the characteristic function exists for
every probability distribution.

%%%%%%%%%%%%%%%%%%%%%%%%%%%%%%%%%%%%%%%%%%%%%%%%%%%%%%%%%%%%
\subsection{Common Moment-Generating Functions}
\label{subsec:common-mgf-table}
%%%%%%%%%%%%%%%%%%%%%%%%%%%%%%%%%%%%%%%%%%%%%%%%%%%%%%%%%%%%

Bernoulli:

\[
\boxed{
M_X(t)
=
1-p+pe^t.
}
\]

Binomial:

\[
\boxed{
M_X(t)
=
(1-p+pe^t)^n.
}
\]

Poisson:

\[
\boxed{
M_X(t)
=
\exp
\left[
\lambda(e^t-1)
\right].
}
\]

Normal:

\[
\boxed{
M_X(t)
=
\exp
\left(
\mu t+\frac12\sigma^2t^2
\right).
}
\]

Exponential with rate \(\lambda\):

\[
\boxed{
M_X(t)
=
\frac{
\lambda
}{
\lambda-t
},
\qquad
t<\lambda.
}
\]

%%%%%%%%%%%%%%%%%%%%%%%%%%%%%%%%%%%%%%%%%%%%%%%%%%%%%%%%%%%%
\subsection{Quick Discrete Distribution Table}
\label{subsec:quick-discrete-distribution-table}
%%%%%%%%%%%%%%%%%%%%%%%%%%%%%%%%%%%%%%%%%%%%%%%%%%%%%%%%%%%%

\begin{center}
\begin{tabular}{c|c|c|c}
\textbf{Distribution}
&
\textbf{Support}
&
\textbf{Mean}
&
\textbf{Variance}
\\ \hline

Bernoulli\((p)\)
&
\(0,1\)
&
\(p\)
&
\(p(1-p)\)
\\

Binomial\((n,p)\)
&
\(0,\ldots,n\)
&
\(np\)
&
\(np(1-p)\)
\\

Geometric\((p)\)
&
\(1,2,\ldots\)
&
\(1/p\)
&
\((1-p)/p^2\)
\\

Poisson\((\lambda)\)
&
\(0,1,\ldots\)
&
\(\lambda\)
&
\(\lambda\)
\end{tabular}
\end{center}

%%%%%%%%%%%%%%%%%%%%%%%%%%%%%%%%%%%%%%%%%%%%%%%%%%%%%%%%%%%%
\subsection{Quick Continuous Distribution Table}
\label{subsec:quick-continuous-distribution-table}
%%%%%%%%%%%%%%%%%%%%%%%%%%%%%%%%%%%%%%%%%%%%%%%%%%%%%%%%%%%%

\begin{center}
\begin{tabular}{c|c|c|c}
\textbf{Distribution}
&
\textbf{Support}
&
\textbf{Mean}
&
\textbf{Variance}
\\ \hline

Uniform\((a,b)\)
&
\([a,b]\)
&
\((a+b)/2\)
&
\((b-a)^2/12\)
\\

Exponential\((\lambda)\)
&
\(x\geq0\)
&
\(1/\lambda\)
&
\(1/\lambda^2\)
\\

Gamma\((\alpha,\beta)\)
&
\(x>0\)
&
\(\alpha/\beta\)
&
\(\alpha/\beta^2\)
\\

Beta\((\alpha,\beta)\)
&
\(0<x<1\)
&
\(\alpha/(\alpha+\beta)\)
&
\(\dfrac{\alpha\beta}{(\alpha+\beta)^2(\alpha+\beta+1)}\)
\\

Normal\((\mu,\sigma^2)\)
&
\(\mathbb{R}\)
&
\(\mu\)
&
\(\sigma^2\)
\end{tabular}
\end{center}

%%%%%%%%%%%%%%%%%%%%%%%%%%%%%%%%%%%%%%%%%%%%%%%%%%%%%%%%%%%%
\subsection{Choosing a Distribution}
\label{subsec:choosing-probability-distribution}
%%%%%%%%%%%%%%%%%%%%%%%%%%%%%%%%%%%%%%%%%%%%%%%%%%%%%%%%%%%%

A useful first-pass guide is:

\begin{itemize}
    \item one binary trial:
    \[
    \boxed{
    \text{Bernoulli}
    }
    \]

    \item number of successes in fixed independent trials:
    \[
    \boxed{
    \text{Binomial}
    }
    \]

    \item trials until first success:
    \[
    \boxed{
    \text{Geometric}
    }
    \]

    \item successes from sampling without replacement:
    \[
    \boxed{
    \text{Hypergeometric}
    }
    \]

    \item counts in a Poisson-process interval:
    \[
    \boxed{
    \text{Poisson}
    }
    \]

    \item waiting time to first Poisson-process event:
    \[
    \boxed{
    \text{Exponential}
    }
    \]

    \item waiting time to several Poisson-process events:
    \[
    \boxed{
    \text{Gamma/Erlang}
    }
    \]

    \item unknown probability on \((0,1)\):
    \[
    \boxed{
    \text{Beta}
    }
    \]

    \item symmetric bell-shaped continuous data:
    \[
    \boxed{
    \text{Normal}
    }
    \]

    \item positive multiplicative data:
    \[
    \boxed{
    \text{Lognormal}
    }
    \]

    \item positive lifetime data:
    \[
    \boxed{
    \text{Weibull}
    }
    \]
\end{itemize}

%%%%%%%%%%%%%%%%%%%%%%%%%%%%%%%%%%%%%%%%%%%%%%%%%%%%%%%%%%%%
\subsection{Common Errors}
\label{subsec:probability-distribution-common-errors}
%%%%%%%%%%%%%%%%%%%%%%%%%%%%%%%%%%%%%%%%%%%%%%%%%%%%%%%%%%%%

\subsubsection{Confusing PMF and PDF}

For a discrete random variable,

\[
\Pr(X=x)
\]

may be positive.

For a continuous random variable,

\[
\boxed{
\Pr(X=x)=0.
}
\]

A density value is not itself a point probability.

%%%%%%%%%%%%%%%%%%%%%%%%%%%%%%%%%%%%%%%%%%%%%%%%%%%%%%%%%%%%
\subsubsection{Ignoring the Support}

A probability formula may only be valid on its stated support.

For example, an exponential random variable cannot take negative values.

%%%%%%%%%%%%%%%%%%%%%%%%%%%%%%%%%%%%%%%%%%%%%%%%%%%%%%%%%%%%
\subsubsection{Confusing Variance and Standard Deviation}

If

\[
X\sim N(\mu,\sigma^2),
\]

then the second parameter is variance under this notation:

\[
\boxed{
\operatorname{Var}(X)=\sigma^2.
}
\]

Standard deviation is \(\sigma\).

%%%%%%%%%%%%%%%%%%%%%%%%%%%%%%%%%%%%%%%%%%%%%%%%%%%%%%%%%%%%
\subsubsection{Confusing Rate and Scale}

For exponential and Gamma distributions, software and textbooks may use rate
or scale.

For the exponential distribution,

\[
\boxed{
\text{scale}
=
\frac1{\text{rate}}.
}
\]

%%%%%%%%%%%%%%%%%%%%%%%%%%%%%%%%%%%%%%%%%%%%%%%%%%%%%%%%%%%%
\subsubsection{Using Binomial for Sampling Without Replacement}

Sampling without replacement creates dependence.

The hypergeometric distribution is generally the exact finite-population
model.

%%%%%%%%%%%%%%%%%%%%%%%%%%%%%%%%%%%%%%%%%%%%%%%%%%%%%%%%%%%%
\subsubsection{Using Geometric Formulas Without Checking the Convention}

Some definitions count trials until success.

Others count failures before success.

This changes the support and mean by \(1\).

%%%%%%%%%%%%%%%%%%%%%%%%%%%%%%%%%%%%%%%%%%%%%%%%%%%%%%%%%%%%
\subsubsection{Assuming Every Distribution Has Finite Mean and Variance}

For example, the Cauchy distribution has neither a finite ordinary mean nor
variance.

%%%%%%%%%%%%%%%%%%%%%%%%%%%%%%%%%%%%%%%%%%%%%%%%%%%%%%%%%%%%
\subsubsection{Treating an Approximation as an Identity}

Normal or Poisson approximations to other distributions are approximations.

Their accuracy depends on the parameter regime.

%%%%%%%%%%%%%%%%%%%%%%%%%%%%%%%%%%%%%%%%%%%%%%%%%%%%%%%%%%%%
\subsection{Exercises}
\label{subsec:probability-distributions-exercises}
%%%%%%%%%%%%%%%%%%%%%%%%%%%%%%%%%%%%%%%%%%%%%%%%%%%%%%%%%%%%

\begin{exercise}[1]
State the PMF of a Bernoulli random variable.
\end{exercise}

\begin{exercise}[1]
State the mean and variance of a Bernoulli random variable.
\end{exercise}

\begin{exercise}[1]
State the PMF of a binomial random variable.
\end{exercise}

\begin{exercise}[1]
Find the mean and variance of

\[
X\sim\operatorname{Binomial}(20,0.3).
\]
\end{exercise}

\begin{exercise}[2]
Find

\[
\Pr(X=4)
\]

for

\[
X\sim\operatorname{Binomial}(10,0.5).
\]
\end{exercise}

\begin{exercise}[2]
Find

\[
\Pr(X\leq2)
\]

for a specified binomial random variable.
\end{exercise}

\begin{exercise}[2]
Find the probability that the first success occurs on trial \(5\) for a
geometric random variable.
\end{exercise}

\begin{exercise}[2]
Verify the memoryless property of the geometric distribution.
\end{exercise}

\begin{exercise}[2]
Compute a hypergeometric probability for a sample drawn without replacement.
\end{exercise}

\begin{exercise}[2]
Find the mean and variance of a Poisson random variable with
\(\lambda=4\).
\end{exercise}

\begin{exercise}[2]
Find

\[
\Pr(X=2)
\]

for

\[
X\sim\operatorname{Poisson}(3).
\]
\end{exercise}

\begin{exercise}[2]
Approximate a binomial probability using a Poisson distribution.
\end{exercise}

\begin{exercise}[2]
Find the mean and variance of

\[
X\sim\operatorname{Uniform}(2,8).
\]
\end{exercise}

\begin{exercise}[2]
Find

\[
\Pr(3<X<6)
\]

for a continuous uniform random variable on \([0,10]\).
\end{exercise}

\begin{exercise}[2]
Find the probability that an exponential waiting time exceeds \(t\).
\end{exercise}

\begin{exercise}[2]
Verify the memoryless property of the exponential distribution.
\end{exercise}

\begin{exercise}[2]
Convert between Gamma rate and scale parameterizations.
\end{exercise}

\begin{exercise}[2]
Find the mean of a Beta\((3,5)\) random variable.
\end{exercise}

\begin{exercise}[2]
Find the variance of a Beta\((3,5)\) random variable.
\end{exercise}

\begin{exercise}[2]
Standardize

\[
X=85
\]

when

\[
X\sim N(70,10^2).
\]
\end{exercise}

\begin{exercise}[2]
Express

\[
\Pr(X\leq x)
\]

in terms of \(\Phi\) for

\[
X\sim N(\mu,\sigma^2).
\]
\end{exercise}

\begin{exercise}[2]
Find the mean of a lognormal random variable from its normal parameters.
\end{exercise}

\begin{exercise}[2]
Show that the exponential distribution is a Gamma distribution with shape
\(1\).
\end{exercise}

\begin{exercise}[3]
Show that a sum of independent Bernoulli variables with common probability
\(p\) is binomial.
\end{exercise}

\begin{exercise}[3]
Show that a sum of independent Poisson variables is Poisson using MGFs.
\end{exercise}

\begin{exercise}[3]
Derive the mean and variance of the binomial distribution.
\end{exercise}

\begin{exercise}[3]
Derive the mean of the geometric distribution.
\end{exercise}

\begin{exercise}[3]
Derive the mean and variance of the Poisson distribution.
\end{exercise}

\begin{exercise}[3]
Verify that the exponential density integrates to \(1\).
\end{exercise}

\begin{exercise}[3]
Verify that the Gamma density integrates to \(1\).
\end{exercise}

\begin{exercise}[3]
Verify that the Beta density integrates to \(1\).
\end{exercise}

\begin{exercise}[3]
Derive the standard normal density from a general normal variable by
standardization.
\end{exercise}

\begin{exercise}[3]
Show that the square of a standard normal variable has a chi-square
distribution with one degree of freedom.
\end{exercise}

\begin{exercise}[3]
Derive the \(t\)-distribution construction from a normal variable and an
independent chi-square variable.
\end{exercise}

\begin{exercise}[3]
Derive the \(F\)-distribution construction from independent chi-square
variables.
\end{exercise}

\begin{exercise}[3]
Find the moment-generating function of a Bernoulli random variable.
\end{exercise}

\begin{exercise}[3]
Find the moment-generating function of a binomial random variable.
\end{exercise}

\begin{exercise}[3]
Find the moment-generating function of a Poisson random variable.
\end{exercise}

\begin{exercise}[3]
Use the normal MGF to derive its mean and variance.
\end{exercise}

\begin{exercise}[4]
Derive the hypergeometric variance including the finite-population correction.
\end{exercise}

\begin{exercise}[4]
Derive the Poisson approximation to the binomial as a limiting result.
\end{exercise}

\begin{exercise}[4]
Study the relation between Gamma waiting times and the Poisson process.
\end{exercise}

\begin{exercise}[4]
Derive the lognormal mean and variance.
\end{exercise}

\begin{exercise}[4]
Derive the Weibull mean using the Gamma function.
\end{exercise}

\begin{exercise}[4]
Study why the Cauchy distribution has no finite ordinary mean.
\end{exercise}

%%%%%%%%%%%%%%%%%%%%%%%%%%%%%%%%%%%%%%%%%%%%%%%%%%%%%%%%%%%%
\subsection{Conceptual Summary}
\label{subsec:probability-distributions-summary}
%%%%%%%%%%%%%%%%%%%%%%%%%%%%%%%%%%%%%%%%%%%%%%%%%%%%%%%%%%%%

Discrete distributions use probability mass functions:

\[
\boxed{
p_X(x)
=
\Pr(X=x).
}
\]

Continuous distributions use probability densities:

\[
\boxed{
\Pr(a\leq X\leq b)
=
\int_a^b f_X(x)\,dx.
}
\]

Important discrete models include:

\[
\boxed{
\text{Bernoulli},
\quad
\text{Binomial},
\quad
\text{Geometric},
\quad
\text{Hypergeometric},
\quad
\text{Poisson}.
}
\]

Important continuous models include:

\[
\boxed{
\text{Uniform},
\quad
\text{Exponential},
\quad
\text{Gamma},
\quad
\text{Beta},
\quad
\text{Normal}.
}
\]

Inferential statistics frequently uses:

\[
\boxed{
\chi^2,
\quad
t,
\quad
F.
}
\]

The standardization

\[
\boxed{
Z
=
\frac{
X-\mu
}{
\sigma
}
}
\]

converts a normal variable into a standard normal variable.

The central practical principle is

\[
\boxed{
\text{match the support, mechanism, and parameterization to the model}.
}
\]

%%%%%%%%%%%%%%%%%%%%%%%%%%%%%%%%%%%%%%%%%%%%%%%%%%%%%%%%%%%%
\subsection{Section Connections}
\label{subsec:probability-distributions-connections}
%%%%%%%%%%%%%%%%%%%%%%%%%%%%%%%%%%%%%%%%%%%%%%%%%%%%%%%%%%%%

\chapterconnection{
Probability Distributions consolidates the standard random-variable models
used throughout Probability, Statistics, Stochastic Processes, Mathematical
Biology, Queueing Theory, Reliability, Finance, and Applied Modeling.
Bernoulli and binomial models connect with discrete probability;
Poisson, exponential, and Gamma distributions connect with counting and
waiting-time processes; normal distributions connect with the Central Limit
Theorem; Beta distributions connect with Bayesian statistics; and chi-square,
Student's \(t\), and \(F\) distributions form the basis of many classical
statistical procedures. The next reference section, Statistical Tables,
collects standard descriptive, inferential, regression, correlation,
confidence-interval, test-statistic, and critical-value formulas.
}

%%%%%%%%%%%%%%%%%%%%%%%%%%%%%%%%%%%%%%%%%%%%%%%%%%%%%%%%%%%%
% SECTION SOURCE TRACKING
%%%%%%%%%%%%%%%%%%%%%%%%%%%%%%%%%%%%%%%%%%%%%%%%%%%%%%%%%%%%

%------------------------------------------------------------
% 14.13 Probability Distributions
%------------------------------------------------------------
%
% Source basis:
%   - Standard probability and mathematical statistics.
%   - Standard discrete and continuous probability distributions.
%
% Used for:
%   - PMFs, PDFs, and CDFs;
%   - expectation and variance;
%   - Bernoulli;
%   - binomial;
%   - geometric;
%   - negative binomial;
%   - hypergeometric;
%   - Poisson;
%   - discrete uniform;
%   - continuous uniform;
%   - exponential;
%   - Gamma and Erlang;
%   - Beta;
%   - normal and standard normal;
%   - lognormal;
%   - Weibull;
%   - Cauchy;
%   - chi-square;
%   - Student's t;
%   - F distribution;
%   - distribution relationships;
%   - MGFs;
%   - characteristic functions;
%   - parameterization warnings;
%   - common errors;
%   - applications and exercises.
%
% Notes:
%   - Statistical Tables is developed next in Section 14.14.
%   - Common Mathematical Theorems follows in Section 14.15.
%
%%%%%%%%%%%%%%%%%%%%%%%%%%%%%%%%%%%%%%%%%%%%%%%%%%%%%%%%%%%%